\section{Introduction}

Visual grounding links a referring expression to the image region it denotes, and it underlies how an assistant connects a user's words to the visible world~\cite{kazemzadeh2014referitgame,liu2024grounding,peng2023kosmos,you2024ferret}. In physical interaction, however, language alone often underdetermines the referent. A request such as \emph{what is that?} contains no description to match, and a cluttered scene may hold several objects that fit any short phrase. People resolve this ambiguity by pointing. Interfaces have combined pointing with speech since \emph{Put-That-There}~\cite{bolt1980put}, and embodied reference understanding now models gesture jointly with language~\cite{chen2021yourefit,li2023touchline}. Pointing-target localisation isolates the spatial core of the problem: given an image that contains a pointing hand, return the region of the indicated object. In egocentric views it is the first step of a gesture-aware assistant, which must establish where the wearer is pointing before it can recognise the object or answer a question about it~\cite{egopoint,choi2026hint,li2026egopointbench}.

Smart glasses are the setting in which this capability is most directly needed. A head-mounted camera shares the wearer's viewpoint, so the pointing hand and the indicated object can be observed in a single frame. The device category is also expanding quickly. IDC reports that worldwide extended-reality shipments grew 44.4\% in 2025, driven primarily by smart glasses, and that glasses without displays already account for the majority of those shipments~\cite{idc2026xr}. Counterpoint Research reports a 263\% year-over-year increase in AI-glasses shipments in the first half of 2026, with display-less models making up 96\% of the total~\cite{counterpoint2026aiglasses}. The growing form factor is therefore a lightweight frame that carries a camera but leaves little room for a battery or for heat dissipation. Most commercial smart glasses carry 200--300\,mAh batteries, less than one twentieth of a typical smartphone's capacity, and sustained computation on them is limited by both battery and thermal budgets~\cite{yuan2026aiglasses}. A full-system power model of an always-on research wearable likewise identifies weight and battery restrictions as the source of stringent power constraints~\cite{lee2025aria2}.

Recent pointing-grounding methods have raised accuracy substantially by building on large vision--language backbones. SV-CoT instantiates visual chain-of-thought inference on Qwen3-VL~\cite{egopoint}; PointVG-R fine-tunes a 7B-parameter Qwen2.5-VL model with chain-of-thought supervision and reinforcement learning~\cite{li2026pointvgr}; and VistaRef adds explicit hand and ray modelling to a BEiT-3 vision--language model whose base variant has 226M parameters~\cite{li2026vistaref}. Without such task-specific training, general-purpose multimodal large language models (MLLMs) remain unreliable on pointing: they often select objects near the hand or visually salient objects instead of the object the gesture indicates~\cite{li2026egopointbench,choi2026hint}. Models at the billion-parameter scale also exceed what current wearable processors can execute~\cite{yuan2026aiglasses}. The 16-bit weights of a 7B-parameter model alone occupy about 14\,GB. Liu et al.~\cite{liu2024mobilellm} note that smartphones offer 6--12\,GB of DRAM, that one application should use no more than about a tenth of it, and that a 7B model is prohibitively expensive on a phone even with 8-bit weights. Glasses have far smaller energy reserves than phones, and current smart-glasses assistants accordingly run lightweight inference locally while delegating multimodal foundation models to the cloud~\cite{yuan2026aiglasses}. Delegating pointing resolution in this way makes the most basic spatial step of the interaction depend on connectivity, adds a network round trip to every query and sends the egocentric frame off the device.

These observations define the question this paper addresses. The pointing-grounding methods above report accuracy for models with hundreds of millions to billions of parameters that also consume a text instruction and, in some settings, an externally supplied hand box~\cite{egopoint}. Sub-million-parameter networks have run on smart-glasses hardware for generic object detection~\cite{moosmann2024tinyissimoyolo}, but generic detection does not decide which object a gesture refers to. We therefore ask \emph{how much pointing-target localisation an image-only model with fewer than one million parameters can provide, and which mechanisms matter at that scale.} Answering it requires more than finding a finger. Egocentric frames can contain hands other than the wearer's~\cite{bambach2015lending}, so the model must anchor its prediction on the wearer's hand. The target box must also capture the referred object's extent, which a fingertip location alone does not determine (Section~\ref{sec:results}). A small model need not replace semantic reasoning. It can pass a target region to a recogniser or assistant while keeping the initial spatial interpretation on the device, a role that is useful only if its accuracy, failure modes and implementation cost are made explicit.

We investigate this question with \model{}, a 722,104-parameter convolutional network that predicts the wearer's hand, fingertip, image-plane direction and target box from a single RGB frame in one forward pass. A shared multi-scale feature map supports the sequence of predictions, and later heads receive features sampled around the selected hand instance rather than only a whole-image descriptor. This structure encodes the dependency of the target on the hand without a separate hand detector or pose estimator. During training, pasted hands act as distractors: the original wearer and target remain the positive reference, and the hand heatmap is penalised for responding to the added hands. Figure~\ref{fig:teaser} shows successful predictions across distinct target configurations.

On the 2,604 target-positive EgoPoint test images, \model{} reaches 59.0$\pm$0.1\% target \pacc{0.5} across three training seeds at $160\times160$ input, with 0.23 counted GFLOPs and 1.6\,MiB of FP16 weights. On prototype glasses, the deployed checkpoint recorded 60--80\,ms inference per frame. Table~\ref{tab:main} places this operating point beside published large models. With roughly $10^4$ times fewer parameters than 7--8B MLLMs, and without a text instruction or a supplied hand box, \model{} falls within the 44.5--63.6\% \pacc{0.5} range reported for general-purpose 7--8B MLLMs used zero-shot or after supervised fine-tuning~\cite{li2026pointvgr}. The compactness has a measurable cost. PointVG-R reports 83.6\% \pacc{0.5}, and the gap widens at the stricter threshold (74.5\% versus 31.5\% \pacc{0.7}). Because the published results use different inputs and evaluation populations, they locate an accuracy regime rather than establish a controlled ranking.

Controlled experiments show that the design choices matter in distinct ways. Removing the instance-conditioned pathway leaves \pacc{0.5} nearly unchanged but reduces target \pacc{0.7} by 2.8 percentage points. Removing distractor training reduces target accuracy by 4.1 points in a synthetic multi-hand stress test. Raising the input to 320 pixels increases \pacc{0.5} to 65.1\% at unchanged parameter count but about four times the counted computation. The evidence also revises an appealing geometric hypothesis: supplying the target head with an explicit ray representation does not improve its localisation after retraining, although removing that branch from an already trained network does affect its predictions. Finally, 69.0\% of target-positive test images place the annotated fingertip \emph{inside} the target box. This is image-plane overlap, not evidence of physical contact, and success on these images must not obscure the lower accuracy when fingertip and target are separated.

Our contributions are:
\begin{itemize}
\item An image-only, single-pass pointing-target localiser for smart glasses that integrates instance-conditioned prediction and wearer-specific distractor supervision within a 722K-parameter CNN, roughly four orders of magnitude smaller than the 7--8B MLLMs reported on this task.
\item A controlled evaluation of localisation precision, input resolution and synthetic hand interference, with three full-model training seeds and four targeted ablations, reported beside published large-model results with their protocol differences stated.
\item An evidence-based account of the design's limits: explicit ray conditioning is not beneficial under the tested recipe, and greater image detail improves localisation without resolving the ambiguity of separated targets.
\end{itemize}
A prototype deployment connects this operating point to smart-glasses inference. We report its observed latency separately from benchmark accuracy and do not infer interaction quality or production battery life from it.
